\documentclass[10pt]{amsart}
\usepackage[a4paper,margin=23mm]{geometry}
\usepackage{amsmath,amssymb,amsthm,hyperref,graphicx,xspace,letltxmacro}
\usepackage[table]{xcolor}
\usepackage[english]{babel}
\providecommand{\th}{}
\newcommand{\ie}{i.e.}\newcommand{\eg}{e.g.}\newcommand{\cf}{cf.}\newcommand{\resp}{resp.}\newcommand{\skpt}{\par\smallskip\noindent}
\newtheorem{theo}{Theorem}[section]
\newtheorem{lemm}[theo]{Lemma}\newtheorem{prop}[theo]{Proposition}\newtheorem{coro}[theo]{Corollary}\newtheorem{conj}[theo]{Conjecture}
\theoremstyle{definition}\newtheorem{rema}[theo]{Remark}\newtheorem{exem}[theo]{Example}
\emergencystretch=3em\pagestyle{plain}
\usepackage{tikz}
\usepackage{array}
\usepackage{multirow}
\usepackage{hhline}
\usepackage{cleveref}
\usepackage{dsfont}

\hfuzz 4.6pt % numerous non visible horizontal overfulls in the tables

\DeclareMathOperator{\conv}{conv}
\DeclareMathOperator{\mon}{m}

\usetikzlibrary{calc,through,backgrounds,shapes,matrix}

\newcolumntype{C}[1]{>{\centering\arraybackslash$}p{#1}<{$}}

\newcommand{\ssm}{\smallsetminus}
\renewcommand{\th}{\textsuperscript{th}}
\newcommand{\ZZ}{\mathbb{Z}}
\newcommand{\RR}{\mathbb{R}}
\newcommand{\Sym}{{\mathfrak{S}}}
\newcommand{\Phiplus}{\Phi^+}
\newcommand{\Phipm}{\Phi_{\geq -1}}
\newcommand{\refl}{\mathcal{R}}

\definecolor{darkblue}{rgb}{0,0,0.7}
\newcommand{\darkblue}{\color{darkblue}}
\definecolor{lightgrey}{rgb}{0.9,0.9,0.9}
\newcommand{\lightgrey}{\cellcolor{lightgrey}}
\definecolor{grey}{rgb}{0.5,0.5,0.5}

\newcommand{\Dfn}[1]{\emph{#1}}

\newcommand{\x}{\mathbf{x}}
\newcommand{\y}{\mathbf{y}}
\renewcommand{\u}{\mathbf{u}}
\newcommand{\f}{\mathbf{f}}
\newcommand{\onevar}{\mathbf{1}}
\newcommand{\id}{\mathds{1}}
\newcommand{\Mpr}{\sfM^{pr}}
\newcommand{\Mex}{\sfM^{ex}}
\newcommand{\Alg}{\mathcal{A}}
\newcommand{\subwordComplex}{\mathcal{SC}}
\newcommand{\clusterComplex}{\subwordComplex\big(\cw{c}\big)}
\newcommand{\greedy}{I_{\operatorname{g}}}
\newcommand{\antigreedy}{I_{\operatorname{ag}}}
\newcommand{\Fpoly}[2]{{\mathsf{F}_{#1}(#2)}}
\newcommand{\FpolyGen}[1]{{\mathsf{F}(#1)}}
\newcommand{\Xvec}[2]{{\mathsf{#1}({#2})}}
\newcommand{\gvec}[1]{\Xvec{g}{#1}}
\newcommand{\dvec}[1]{\Xvec{d}{#1}}
\newcommand{\cvec}[1]{\Xvec{c}{#1}}
\newcommand{\zetag}{{\zeta_{\operatorname{g}}}}
\newcommand{\zetaag}{{\zeta_{\operatorname{ag}}}}
\newcommand{\sq}[1]{{\rm #1}}
\newcommand{\Q}{\sq{Q}}
\newcommand{\q}{\sq{q}}
\newcommand{\s}{\sq{s}}
%\newcommand{\w}{\sq{w}}
\newcommand{\sref}{\mathcal{S}}
\newcommand{\sfr}{\mathsf{r}}
\newcommand{\sfR}{\mathsf{R}}
\newcommand{\sfw}{\mathsf{w}}
\newcommand{\sfW}{\mathsf{W}}
\newcommand{\sfB}{\mathsf{B}}
\newcommand{\wo}{{w_\circ}}
\renewcommand{\c}{\sq{c}}
\newcommand{\sw}[2]{\sq{#1}(\sq{#2})}
\newcommand{\cwo}[1]{\sw{\wo}{#1}}
\newcommand{\cw}[1]{\sq{#1}\cwo{#1}}
\newcommand{\brickVector}{\sfB}
\newcommand{\brickPolytope}{\mathcal{B}}
\newcommand{\Root}[2]{{\sfr}(#1,#2)}
\newcommand{\Roots}[1]{{\sfR}(#1)}
\newcommand{\Weight}[2]{{\sfw}(#1,#2)}
\newcommand{\Weights}[1]{{\sfW}(#1)}
\newcommand{\IntervalWeight}[1]{{\operatorname{Int}\big(\Weight{\greedy}{#1},\ \Weight{\antigreedy}{#1}\big)}}
\newcommand{\coRoot}[2]{{\sfr^\vee}(#1,#2)}
\newcommand{\coRoots}[1]{{\sfR^\vee}(#1)}
\newcommand{\coWeight}[2]{{\sfw^\vee}(#1,#2)}
\newcommand{\coWeights}[1]{{\sfW^\vee}(#1)}
\newcommand{\Newton}[1]{{\operatorname{Newton}\big(#1 \big)}}
\newcommand{\set}[2]{\left\{ #1 \;:\; #2 \right\}}
\newcommand{\bigset}[2]{\big\{ #1 \;:\; #2 \big\}}
\newcommand{\bigmultiset}[2]{\big\{\!\!\big\{ #1 \;:\; #2 \big\}\!\!\big\}}
\newcommand{\Mink}{{\sum}}
\newcommand{\dotprod}[2]{\left\langle \, #1 \,\middle|\, #2 \, \right\rangle}
\newcommand{\bigdotprod}[2]{\big\langle \, #1 \,\big|\, #2 \, \big\rangle}
\newcommand{\wordprod}[2]{\Pi{#1}_{#2}}
\newcommand{\sfM}{\mathsf{M}}
\newcommand{\mutmatrix}{\widetilde\sfM}


\Crefname{conj}{Conjecture}{Conjectures}
\Crefname{theo}{Theorem}{Theorems}
\Crefname{rema}{Remark}{Remarks}
\Crefname{coro}{Corollary}{Corollaries}
\Crefname{lemm}{Lemma}{Lemmas}
\Crefname{observation}{Observation}{Observations}
\Crefname{hope}{Hope}{Hopes}
\Crefname{prop}{Proposition}{Propositions}
\Crefname{exam}{Example}{Examples}


\graphicspath{{./figures/}}

\newcommand*{\mk}{\mkern -1mu}
\newcommand*{\Mk}{\mkern -2mu}
\newcommand*{\mK}{\mkern 1mu}
\newcommand*{\MK}{\mkern 2mu}

\newcommand{\SourcePx}[1]{\SourcePxAux#1\relax}
\def\SourcePxAux#1px\relax{#1pt}
\newcolumntype{C}[1]{>{\centering\arraybackslash$}p{\SourcePx{#1}}<{$}}
\makeatletter
\expandafter\def\csname sourceCref@ex:subwordA2\endcsname{Example 2.24}
\expandafter\def\csname sourceCref@sec:proofextendedpart\endcsname{Section 3}
\expandafter\def\csname sourceCref@fig:A4trianexample\endcsname{Figure 1}
\expandafter\def\csname sourceCref@fig:132Tpaths\endcsname{Figure 2}
\expandafter\def\csname sourceCref@ex:A3two\endcsname{Example 4.8}
\expandafter\def\csname sourceCref@ex:shortestpaths\endcsname{Example 4.13}
\expandafter\def\csname sourceCref@weightsc132A3\endcsname{Figure 3}
\expandafter\def\csname typedCref@lem:commutation\endcsname{Lemma~\ref{lem:commutation}}
\expandafter\def\csname typedCref@thm:minkowski\endcsname{Theorem~\ref{thm:minkowski}}
\expandafter\def\csname typedCref@thm:cluster1\endcsname{Theorem~\ref{thm:cluster1}}
\expandafter\def\csname typedCref@thm:cluster2\endcsname{Theorem~\ref{thm:cluster2}}
\expandafter\def\csname typedCref@thm:exchangematrix\endcsname{Theorem~\ref{thm:exchangematrix}}
\expandafter\def\csname typedCref@thm:extendedpart\endcsname{Theorem~\ref{thm:extendedpart}}
\expandafter\def\csname typedCref@cor:extendedpart\endcsname{Corollary~\ref{cor:extendedpart}}
\expandafter\def\csname typedCref@conj:newtonpolytope\endcsname{Conjecture~\ref{conj:newtonpolytope}}
\expandafter\def\csname typedCref@conj:latticepoints\endcsname{Conjecture~\ref{conj:latticepoints}}
\expandafter\def\csname typedCref@thm:main1\endcsname{Theorem~\ref{thm:main1}}
\expandafter\def\csname typedCref@thm:main2\endcsname{Theorem~\ref{thm:main2}}
\expandafter\def\csname typedCref@cor:dvector\endcsname{Corollary~\ref{cor:dvector}}
\expandafter\def\csname typedCref@cor:gvector\endcsname{Corollary~\ref{cor:gvector}}
\expandafter\def\csname typedCref@prop:clustervariables\endcsname{Proposition~\ref{prop:clustervariables}}
\expandafter\def\csname typedCref@thm:newton\endcsname{Theorem~\ref{thm:newton}}
\expandafter\def\csname typedCref@lem:lemma33\endcsname{Lemma~\ref{lem:lemma33}}
\expandafter\def\csname typedCref@prop:extendedpartinitial\endcsname{Proposition~\ref{prop:extendedpartinitial}}
\expandafter\def\csname typedCref@prop:exchangeRelationslemma\endcsname{Proposition~\ref{prop:exchangeRelationslemma}}
\expandafter\def\csname typedCref@lem:weightequality\endcsname{Lemma~\ref{lem:weightequality}}
\expandafter\def\csname typedCref@lem:positiveweightshift\endcsname{Lemma~\ref{lem:positiveweightshift}}
\expandafter\def\csname typedCref@lem:rootlemma\endcsname{Lemma~\ref{lem:rootlemma}}
\expandafter\def\csname typedCref@thm:schiffler\endcsname{Theorem~\ref{thm:schiffler}}
\expandafter\def\csname typedCref@lem:diagonalspositiveroots\endcsname{Lemma~\ref{lem:diagonalspositiveroots}}
\expandafter\def\csname typedCref@cor:TpathAn\endcsname{Corollary~\ref{cor:TpathAn}}
\expandafter\def\csname typedCref@prop:greedytpath\endcsname{Proposition~\ref{prop:greedytpath}}
\expandafter\def\csname typedCref@cor:TpathAn2\endcsname{Corollary~\ref{cor:TpathAn2}}
\expandafter\def\csname typedCref@lem:greedyantigreedyweight\endcsname{Lemma~\ref{lem:greedyantigreedyweight}}
\expandafter\def\csname typedCref@prop:weightsAn\endcsname{Proposition~\ref{prop:weightsAn}}
\expandafter\def\csname typedCref@cor:intervalequality\endcsname{Theorem~\ref{cor:intervalequality}}
\expandafter\def\csname typedCref@conj:newtonpolytope,conj:latticepoints\endcsname{Conjectures~\ref{conj:newtonpolytope} and~\ref{conj:latticepoints}}
\LetLtxMacro\NativeCref\Cref
\LetLtxMacro\NativePageref\pageref
\newcommand{\SourceCref}[1]{\ifcsname sourceCref@#1\endcsname\href{https://doi.org/10.5802/alco.25}{\csname sourceCref@#1\endcsname\ of the source article}\else\ifcsname typedCref@#1\endcsname\csname typedCref@#1\endcsname\else\NativeCref{#1}\fi\fi}
\newcommand{\SourcePageref}[1]{\ifcsname sourceCref@#1\endcsname\href{https://doi.org/10.5802/alco.25}{557 of the source article}\else\NativePageref{#1}\fi}
\AtBeginDocument{\LetLtxMacro\NativeCref\Cref\LetLtxMacro\NativePageref\pageref\LetLtxMacro\Cref\SourceCref\LetLtxMacro\pageref\SourcePageref}
\makeatother
\begin{document}
\title{Newton polytopes and saturated monomial supports of principal-coefficient type-A cluster F-polynomials from subword weights}
\author{Sarah B. Brodsky \and Christian Stump}
\maketitle
\setcounter{section}{1}
\section{Coxeter roots, principal coefficients and cluster subwords}
\subsection{Root systems}\label{sec:rootsystem}

Let~$(W,\sref)$ be a finite crystallographic Coxeter system acting essentially on a Euclidean vector space~$V$ of dimension~$n$ with inner product~$\dotprod{\cdot}{\cdot}$, we refer to~\cite[Part~II]{Hum1990} for definitions of Coxeter systems and their geometric representations. Let $\Delta = \set{\alpha_s}{s \in \sref}$ be a choice of \Dfn{simple roots} and set~$\Delta^\vee = \set{\alpha_s^\vee}{s \in \sref}$ to be the \Dfn{simple coroots}. We then have that $\alpha_s^\vee = 2\alpha_s / \dotprod{\alpha_s}{\alpha_s}$, that the \Dfn{Cartan matrix} $A = (a_{st})_{s,t \in \sref}$ is given by $a_{st} = \dotprod{\alpha_t}{\alpha_s^\vee}$, and that $s(\alpha_t) = \alpha_t - a_{st} \alpha_s$. The \Dfn{fundamental weights} $\nabla = \set{\omega_s}{s \in \sref}$ and \Dfn{fundamental coweights} $\nabla^\vee = \set{\omega_s^\vee}{s \in \sref}$ are the bases dual to the simple coroots and to the simple roots, respectively. That is, $\dotprod{\omega_s}{\alpha_t^\vee} = \dotprod{\alpha_s}{\omega_t^\vee} = \delta_{s=t}$. It is then easy to verify that
\begin{align}
\alpha_s = \sum_{t \in \sref}a_{ts}\omega_t, \quad \alpha_t^\vee = \sum_{s \in \sref}a_{ts}\omega_s^\vee, \label{eq:rootweight}
\end{align}
and that moreover, $s(\omega_t) = \omega_t - \delta_{s=t}\alpha_s$ for $s \in \sref$.

Let~$c \in W$ be a \Dfn{(standard) Coxeter element} for $(W,\sref)$; i.e.~$c = s_1\dots s_n$ is the product of all elements in~$\sref$ in some order, and fix the reduced word $\c = \s_1\dots\s_n$ of~$c$. We then often write $\alpha_i = \alpha_{s_i}$, $\alpha_i^\vee = \alpha_{s_i}^\vee$, $\omega_i = \omega_{s_i}$, and $\omega_i^\vee = \omega_{s_i}^\vee$ to avoid double indices. Denote by $L = \ZZ\Delta$ the $\ZZ$-lattice spanned by $\Delta$, by~$L^+$ the nonnegative span $\ZZ_{\geq 0}\Delta$, and by~$L^- = -L^+$ the nonpositive span. We call $\beta \in L$ \Dfn{sign-coherent} if $\beta \in L^+ \sqcup L^-$. Denote moreover by $\Phi = W(\Delta) = \bigset{w(\alpha_s)}{w \in W, s \in \sref} \subseteq L^+ \sqcup L^-$ the \Dfn{root system} for $(W,\sref)$, by $\Phiplus = \Phi \cap L^+$ the \Dfn{positive roots}, and by $\Phipm = \Phiplus \sqcup -\Delta \subseteq \Phi$ the \Dfn{almost positive roots}. The set of reflections in~$W$ is denoted by
\[
\refl = \set{ wsw^{-1} }{ w \in W, s \in \sref}.
\]
The map from $\Phi$ to $\refl$ given by sending $\beta = w(\alpha_s) \in \Phi$ to $s_\beta = w s w^{-1} \in \refl$ is two-to-one and restricts to a bijection between~$\Phiplus$ and~$\refl$. We often use the letter $N = |\Phiplus| = | \refl |$, so that $n+N = |\Phipm|$.
\subsection{Coxeter elements and skew-symmetrizable Cartan matrices}\label{sec:coxeterelement}

For simple reflections $s \neq t$, we have that $a_{st} = 0$ if and only if $st = ts$. We can thus think of the Coxeter element~$c$ as an acyclic orientation of the corresponding Coxeter diagram by orienting an edge $s \rightarrow t$ if~$s$ comes before~$t$ in the given reduced word~$\c = \s_1 \s_2 \dots \s_n$ (or, equivalently, in all reduced words for~$c$). It is indeed the case that this mapping yields a one-to-one correspondence between Coxeter elements and acyclic orientations of the Coxeter diagram, and that the reduced words for a Coxeter element are given by all linear extensions of this orientation. The reason we work with the particular word $\c = \s_1\dots\s_n$ is that it will later simplify several notations, both in the indexing of variables, and in the definition of subword complexes (see \Cref{sec:subwordcomplexes}).

For a given such orientation of the Coxeter diagram, define the skew-symmetri\-zable matrix $\sfM_c = (b_{st})_{s,t \in \sref}$ by
\begin{align}
b_{st} =
\begin{cases}
-a_{st} & \text{if } s \rightarrow t,\\
\phantom{-}a_{st} & \text{if } s \leftarrow t,\\
\phantom{-}0 & \text{else}.
\end{cases} \label{eq:bmatrix}
\end{align}

\subsection{Cluster algebras and cluster complexes}\label{sec:cluster}

To a skew-symmetrizable Cartan matrix $\sfM_c$, one may associate an \Dfn{initial (cluster) seed} $\big(\mutmatrix_c, \x, \y \big)$, where $\mutmatrix_c$ is the \Dfn{exchange matrix} $\big[\begin{smallmatrix} \sfM_c \\ \id_n \end{smallmatrix}\big]$ with principal part~$\sfM_c$ and extended part~$\id_n$ an identity matrix, and where $\x = (x_1,\dots,x_n)$ are the cluster variables (the cluster of the seed), and $\y = (y_1,\dots,y_n)$ are the frozen variables (the coefficients of the seed).
Cluster seeds may now be mutated in directions~$1,\dots,n$, we refer to \Cref{sec:proofextendedpart} for the technical definitions of cluster mutations. The \Dfn{cluster algebra} $\Alg(W,c) = \Alg(\mutmatrix_c)$ is then the subalgebra of the ring of rational functions generated by all cluster variables obtained from the cluster variables in the initial seed by mutations, we refer to~\cite{FZ20022,FZ2003,FZ2007} for all needed background on cluster algebras. In the present context, one should think of the variables~$x_k$ and~$y_k$ as being indexed by~$\alpha_k$ for $1 \leq k \leq n$, so they are in particular indexed in a way that is consistent with the order of the simple reflections in the given Coxeter element~$c$.

%\medskip
It is known that every cluster variable~$u(\x,\y) \in \Alg(W,c)$ is a Laurent polynomial in the ring $\ZZ[x_1^{\pm 1},\dots,x_n^{\pm 1};y_1,\dots,y_n]$, i.e. $u(\x,\y) = p(\x,\y)/\mon(\x)$ where~$p(\x,\y)$ is a polynomial in~$\x,\y$ with integer coefficients and $\mon(\x)$ is a monomial in~$\x$, see~\cite[Proposition~3.6]{FZ2007}. The \Dfn{$d$-vector} $\dvec{u}$ of~$u(\x,\y)$ is the exponent vector of the denominator monomial $\mon(\x)$, i.e. $\dvec{u} = (d_1,\dots,d_n)$ for $\mon(\x) = x_1^{d_1}\dots x_n^{d_n}$ and should be thought of as a vector in the basis $\Delta$, i.e.
\[
\dvec{u} = d_1\alpha_{1} + \dots + d_n\alpha_{n}.
\]
Under this identification, it is shown in~\cite[Theorem~1.9]{FZ2003} and~\cite{YZ2008}, that the map $u \mapsto \dvec{u}$ is a bijection between cluster variables in $\Alg(W,c)$ and almost positive roots~$\Phipm$ where one sets $\dvec{u} = -\alpha_i \in -\Delta$ for $u(\x,\y) = x_i$. We will regularly use this bijection in indexing objects. For example, we denote by
\[
\Fpoly{u}{\y} = \Fpoly{\beta}{\y} = u(\onevar,\y) = p(\onevar,\y)
\]
the \Dfn{$F$-polynomial} associated to~$u(\x,\y) \in \Alg(W,c)$ and to $\beta \in \Phipm$ with $\dvec{u} = \beta$. As $\Fpoly{\beta}{\y} = 1$ for $\beta \in -\Delta$, one usually considers $F$-polynomials only for positive roots $\beta \in \Phiplus$. We also denote by $\gvec{u} = \gvec{\beta} = (g_1,\dots,g_n)$ the \Dfn{$g$-vector} given by the exponent vector of~$u(\x,0)$. For reasons that will become clear later, we consider the $g$-vector to live inside the weight space, i.e. $\gvec{u} = g_1\omega_1+\dots+g_n\omega_n$.
\subsection{Subword complexes}\label{sec:subwordcomplexes}

Let~$\Q = \q_1\dots\q_m$ be a word in the simple system~$\sref$ and let $\rho \in W$. The \Dfn{subword complex} $\subwordComplex(\Q,\rho)$ is the simplicial complex of (positions of) letters in~$\Q$ whose complement contains a reduced word of~$\rho$. These complexes were introduced by A.~Knutson and E.~Miller in~\cite{KM2004}, we refer to \Cref{ex:subwordA2} on page~\pageref{ex:subwordA2} for a detailed example for~$\rho = \wo \in W$ being the longest element. Observe that $\subwordComplex(\Q,\rho)$ is pure by construction with ground set $[m] = \{1,\dots,m\}$ given by the indices of letters in~$\Q$. Its facets thus all have the same size and we consider them as sorted lists of integers, written in set notation. This is, $I = \{i_1 < \dots < i_n\}$ is a facet of $\subwordComplex(\Q,\rho)$ if and only if the word $\q_1\dots \widehat \q_{i_1} \dots \widehat \q_{i_n} \dots \q_m$, with the letters $\q_{i_1},\dots,\q_{i_n}$ omitted, is a reduced word for~$\rho \in W$.

%\medskip

Recall the following fundamental observation about subword complexes. It explains in all later constructions the independence of the chosen reduced word~$\c = \s_1\dots\s_n$ of the Coxeter element~$c$.

\begin{lemm}[{\cite[Proposition~3.8]{CLS2014}}]\label{lem:commutation}
Let~$\Q'$ be a word in~$\sref$ that coincides with~$\Q$ up to commutations of consecutive letters that commute in~$W$. Then $\subwordComplex(\Q,\rho) \cong \subwordComplex(\Q',\rho)$, and the isomorphism is given by the natural identification between letters in~$\Q$ and in~$\Q'$.
\end{lemm}

In this paper, we are only interested in the case that~$\rho = \wo \in W$ is the unique longest element with respect to the weak order, and $\Q$ being one specific word constructed from the Coxeter element~$c$. We thus write $\subwordComplex(\Q)$ for $\subwordComplex(\Q,\wo)$ and assume that~$\Q$ does indeed contain a reduced word for~$\wo$. This immediately implies that $\subwordComplex(\Q)$ is a simplicial sphere, see~\cite[Theorem~3.7]{KM2004}. Define~$\greedy$ and $\antigreedy$ to be the lexicographically first and last facets of $\subwordComplex(\Q)$, respectively. These are called the \Dfn{greedy facet} and the \Dfn{antigreedy facet}.

For $\Q = \q_1 \dots \q_m$, associate to any facet~$I$ of the subword complex~$\subwordComplex(\Q)$ a \Dfn{root function} $\Root{I}{\cdot\,} : [m] \to W(\Delta)$ and a \Dfn{weight function}~$\Weight{I}{\cdot\,} : [m] \to W(\nabla)$ defined by
\[
\Root{I}{k} = \wordprod{\Q}{[k-1] \ssm I}(\alpha_{q_k}) \quad \text{and} \quad \Weight{I}{k} = \wordprod{\Q}{[k-1] \ssm I}(\omega_{q_k}),
\]
where~$\wordprod{\Q}{X}$ denotes the product of the simple reflections~$q_x \in \Q$, for~$x \in X$, in the order given by~$\Q$. For later convenience, we as well define the \Dfn{coroot function} $\coRoot{I}{\cdot\,} : [m] \to W(\Delta^\vee)$ and a \Dfn{coweight function}~$\coWeight{I}{\cdot\,} : [m] \to W(\nabla^\vee)$ by
\[
\coRoot{I}{k} = \wordprod{\Q}{[k-1] \ssm I}(\alpha_{q_k}^\vee) \quad \text{and} \quad \coWeight{I}{k} = \wordprod{\Q}{[k-1] \ssm I}(\omega_{q_k}^\vee).
\]

Observe that it is immediate from this definition that all of these functions are invariant under the isomorphism given in \Cref{lem:commutation}. The root function (and, equivalently, the coroot function) locally encodes the flip property in the subword complex: each facet adjacent to~$I$ in~$\subwordComplex(\Q)$ is obtained by exchanging an element~${i \in I}$ with the unique element~$j \notin I$ such that $\Root{I}{j} \in \{ \pm \Root{I}{i} \}$. If $i < j$ such a flip is called \Dfn{increasing}, and it is called \Dfn{decreasing} otherwise. Observe that the greedy facet and the antigreedy facet are the unique facets such that every flip is increasing and decreasing, respectively. After this exchange, the root function and the weight function are updated by an application of~$s_{\Root{I}{i}}$ as recalled in \Cref{lem:lemma33} below. The root and the coroot functions are used to define the \Dfn{root configuration} and the \Dfn{coroot configuration} of the facet~$I$ as the ordered multisets
\begin{align}
\Roots{I} = \bigmultiset{\Root{I}{i}}{i \in I},\quad \coRoots{I} = \bigmultiset{\coRoot{I}{i}}{i \in I}. \label{eq:rootconfiguration}
\end{align}

Similarly, the weight and the coweight functions are used to define the \Dfn{weight configuration} and the \Dfn{coweight configuration}
\begin{align}
\Weights{I} = \bigmultiset{\Weight{I}{i}}{i \in I},\quad \coWeights{I} = \bigmultiset{\coWeight{I}{i}}{i \in I}. \label{eq:weightconfiguration}
\end{align}

By ordered multiset we simply mean the ordered tuple written in set notation. For later convenience, we denote by
\[
\Root{I}{i}_j = \dotprod{\Root{I}{i}}{\omega_j^\vee}
\]
the coefficient of~$\alpha_j$ in the root $\Root{I}{i}$.
\subsection{Cluster complexes as subword complexes}\label{sec:clustercomplexes}

For the Coxeter element~$c$ with fixed reduced word $\c = \s_1\dots\s_n$, the \Dfn{Coxeter-sorting word} (or \Dfn{$\c$-sorting word}) $\rho(\c)$ of an element $\rho \in W$ is given by the lexicographically first subword of $\c^\infty$ that is a reduced word for~$\rho$. In particular, $\cwo{c}$ is the lexicographically first subword of $\c^\infty$ that is a reduced word for the longest element~$\wo \in W$. Observe that the word $\rho(\c)$ does also depend on the chosen reduced word. But instead, one should think of this sorting word as being associated to the Coxeter element~$c$ and being defined up to commutations of consecutive commuting letters. The notion of $c$-sorting words was defined by N.~Reading in~\cite{R2007} and plays a crucial role in the combinatorial descriptions of finite type cluster algebras and in particular in the description of cluster complexes in terms of subword complexes. The main results in~\cite{CLS2014} provides the following description of the combinatorics of the cluster complex of $\Alg(W,c)$, where we observe from \Cref{lem:commutation} that the subword complex does not depend on the chosen reduced word for~$c$.

\begin{theo}[{\cite[Theorem~2.2]{CLS2014}}]\label{thm:cluster1}
The cluster complex of the cluster algebra $\Alg(W,c)$ is isomorphic to the subword complex $\clusterComplex$.
\end{theo}

We thus refer to $\clusterComplex$ as the \Dfn{$c$-cluster complex}. We moreover remark that the abstract simplicial complex $\clusterComplex$ does not depend on the chosen Coxeter element~$c$, while its combinatorics in the sense of root and weight functions does depend on~$c$. This is, $\subwordComplex\big(\cw{c}\big) \cong \subwordComplex\big(\cw{c'}\big)$ for any two Coxeter elements~$c$ and~$c'$ with reduced words~$\c$ and~$\c'$, respectively, see~\cite[Theorem~2.6]{CLS2014}.

%\medskip

One identifies positions in $\cw{c}$ and almost positive roots by sending the $k$\th\ letter $\s_{k}$ ($1 \leq k \leq n$) of the initial copy of~$\c = \s_1\dots\s_n$ to the negative simple root $-\alpha_k = -\alpha_{s_k}$, and the $k$\th\ letter $\q_k$ ($1 \leq k \leq N$) of $\cwo{c} = \q_1\dots\q_N$ to the positive root $q_1\dots q_{k-1}(\alpha_{q_k})$. See \Cref{lem:lemma33}$\MK$\eqref{it:bijectionpositiveroots} on page~\pageref{lem:lemma33} that this indeed is a bijection, and observe that this equals in the natural way the root function of the greedy facet. In symbols, this is for $1 \leq k \leq N$
\begin{align}
q_1\dots q_{k-1}(\alpha_{q_k}) = \Root{\greedy}{n+k}. \label{eq:rootbijection}
\end{align}

That identification yields the isomorphism in \Cref{thm:cluster1} by sending a cluster to the positions inside the word~$\cw{c}$ corresponding to the almost positive roots of the $d$-vectors of the cluster. To make this explicit, we use the following notation. Let~$I = \{ i_1 < \dots < i_n\}$ be a facet of the cluster complex $\clusterComplex$. We then denote by $S(I) = \big(\mutmatrix(I),\u(I),\f(I)\big)$ with
\begin{align*}
\mutmatrix(I) &=
\begin{bmatrix} \Mpr(I) \\ \Mex(I)
\end{bmatrix} \\
\u(I) &= \big(u_{i_1}(I),\dots,u_{i_n}(I)\big) \\
\f(I) &= \big(f_{i_1}(I),\dots,f_{i_n}(I)\big)
\end{align*}
the cluster seed of $\Alg(W,c)$ corresponding to~$I$ under the given isomorphism between cluster variables, almost positive roots, and positions in the word $\cw{c}$. The columns of $\mutmatrix(I)$ are then also indexed by the positions $i_1,\dots,i_n$ of~$I$ as are the rows of $\Mpr(I)$, while the rows of $\Mex(I)$ are indexed by the positions $1,\dots,n$ (which are the positions of the greedy facet~$\greedy = \{1,\dots,n\}$). We also denote by~$\cvec{I,i}$ the $c$-vector coming from column $i \in I$ of $\Mex(I)$, and by $\gvec{I,i}$ the $g$-vector of the entry $u_i(I)$.
To state the main conjecture of this paper, we define the \Dfn{Newton polytope} of an $F$-polynomial $\Fpoly{\beta}{\y}$ as the convex hull of its exponent vectors in the root basis~$\Delta$. That is,
\[
\Newton{\Fpoly{\beta}{\y}} = \conv\bigset{\lambda_1 \alpha_1+\dots+\lambda_n\alpha_n}{y_1^{\lambda_1}\dots y_n^{\lambda_n}\text{ monomial in }\Fpoly{\beta}{\y}}.
\]
\setcounter{theo}{11}
\begin{conj}\label{conj:newtonpolytope}
Let $\Fpoly{\beta}{\y}$ be the $F$-polynomial associated to the positive root~$\beta$ for the cluster algebra $\Alg(W,c)$. Let~$k$ be the unique index $k \in \{n+1,\dots,n+N\}$ such that $\Root{\greedy}{k} = \beta$ associated to~$\beta$ in Equation~\eqref{eq:rootbijection}. Then
\[
\Newton{\Fpoly{\beta}{\y}} = \conv\set{\Weight{I}{k} - \Weight{\antigreedy}{k}}{I \text{ facet of }\clusterComplex}.
\]
\end{conj}
\begin{conj}\label{conj:latticepoints}
The exponent vectors of the monomials of the $F$-polynomial $\Fpoly{\beta}{\y}$ are given by all lattice points inside its Newton polytope,
\[
\bigset{\lambda_1 \alpha_1+\dots+\lambda_n\alpha_n}{y_1^{\lambda_1}\dots y_n^{\lambda_n}\text{ monomial in }\Fpoly{\beta}{\y}} = \Newton{\Fpoly{\beta}{\y}} \cap L^+.
\]
\end{conj}
\begin{theo}\label{thm:main1}
\Cref{conj:newtonpolytope,conj:latticepoints} hold for~$\Alg(W,c)$ with~$W$ of type~$A_n$.
\end{theo}
\section{Required root and weight core}
\begin{lemm}[{\cite[Lemma~3.3 \& Lemma~3.6]{CLS2014}, \cite[Lemma~3.3, Lemma~4.4 \& Proposition~6.6]{PS2015}}]\label{lem:lemma33}
Let~$I$ and~$J$ be two adjacent facets of the subword complex $\subwordComplex(\Q)$ with $I\setminus i = J \setminus j$. Then
\begin{enumerate}
\item\label{it:bijectionpositiveroots}
The map $\Root{I}{\cdot\,}: k \mapsto \Root{I}{k}$ is a bijection between the complement of~$I$ and~$\Phi^+$.
\item\label{it:rootflip}
The position~$j$ is the unique position in the complement of~$I$ for which $\Root{I}{j}\in \{\pm \Root{I}{i}\}$. Moreover, $\Root{I}{j}=\Root{I}{i}\in \Phi^+$ if $i<j$, while $\Root{I}{j}=-\Root{I}{i}\in \Phi^-$ if $j<i$.
\item\label{it:rootupdate}
The map $\Root{J}{\cdot\,}$ is obtained from $\Root{I}{\cdot\,}$ by
\[
\Root{J}{k}=\left\{
\begin{array}{cc}s_{\Root{I}{i}}(\Root{I}{k}) & \mbox{if } \min\{i,j\}<k\leq \max\{i,j\}, \\\Root{I}{k} & \mbox{otherwise}.
\end{array}\right.
\]
\item\label{it:weightupdate}
The map~$\Weight{J}{\cdot\,}$ is obtained from~$\Weight{I}{\cdot\,}$~by
\[
\Weight{J}{k} =
\begin{cases} s_{\Root{I}{i}}(\Weight{I}{k}) & \text{if } \min\{i,j\} < k \le \max\{i,j\}, \\ \Weight{I}{k} & \text{otherwise}.
\end{cases}
\]
\item\label{it:weightroot}
For $k \in I$, we have for $k' \in I$ with $k' \neq k$ that
\[
\dotprod{\Root{I}{k'}}{\Weight{I}{k}} = 0
\]
and we have for $k' \notin I$ that
\[
\begin{cases}
\dotprod{\Root{I}{k'}}{\Weight{I}{k}} \geq 0 & \text{ if }k' \ge k, \\
\dotprod{\Root{I}{k'}}{\Weight{I}{k}} \leq 0 & \text{ if }k' < k.
\end{cases}
\]
\end{enumerate}
\end{lemm}
After this calculation, we present several general lemmas about cluster complexes that we then use in \Cref{sec:proofmain1} in type~$A_n$ to deduce \Cref{thm:main1}.

\begin{lemm}\label{lem:weightequality}
Let $I,J$ be two facets of $\clusterComplex$ with $k \in I \cap J$. Then $\Weight{I}{k} = \Weight{J}{k}$.
\end{lemm}

\begin{proof}
This is a direct consequence of \Cref{lem:lemma33}$\MK$\eqref{it:weightupdate} and~\eqref{it:weightroot} and the observation that all facets of $\clusterComplex$ containing~$k$ are connected by flips (see~\cite[Corollary~3.11]{PS2015}). (Indeed, this property of the weight function was already used in the proof of~\cite[Proposition~6.8]{PS2015}.) To see this, assume that $k \in I \setminus i = J \setminus j$. The first part of \Cref{lem:lemma33}$\MK$\eqref{it:weightupdate} shows how the weight $\Weight{J}{k}$ is obtained from $\Weight{I}{k}$ by applying $s_{\Root{I}{i}}$ and the first part of~\eqref{it:weightroot} implies that $\Weight{I}{k}$ is contained in the hyperplane fixed by $s_{\Root{I}{i}}$, implying
\[
\Weight{J}{k} = s_{\Root{I}{i}}(\Weight{I}{k}) = \Weight{I}{k}. \qedhere
\]
\end{proof}

This lemma implies as well that $\Weight{I}{i} = \Weight{\antigreedy}{i}$ for any facet~$I$ and any $i \in I$. If $i \in \antigreedy$, this follows immediate. Otherwise, this follows from the observation that one can construct a facet $I'$ with $i \in I'$ being minimal. \Cref{lem:weightequality} then implies that $\Weight{I}{i} = \Weight{I'}{i}$ and \Cref{lem:lemma33}$\MK$\eqref{it:weightupdate} implies that $\Weight{I'}{i} = \Weight{\antigreedy}{i}$. We do not make further use of this additional information.

%\medskip

We next recall the following lemma.

\begin{lemm}[{\cite[Lemma~4.5]{PS2015}}]\label{lem:positiveweightshift}
Let $I \setminus i=J\setminus j$ with $i < j$ be two facets of $\clusterComplex$. For any $k \in \{1,\dots,n+N\}$ we then have
\[
\Weight{J}{k} = \Weight{I}{k} - \lambda\Root{I}{i} \text{ for some } \lambda\in\ZZ_{\geq 0} \text{ and } \Root{I}{i} \in \Phiplus.
\]
In particular, $\Weight{I}{k} - \Weight{J}{k} \in L^+$.
\end{lemm}

\begin{proof}
This is a direct consequence of \Cref{lem:lemma33}$\MK$\eqref{it:weightupdate} and~\eqref{it:weightroot}.
\end{proof}

The following lemma has not been considered before and will serve as the starting point of understanding $F$-polynomials in terms of the weight function.

\begin{lemm}\label{lem:rootlemma}
For $k \in \{n+1,\dots,n+N\}$, we have that
\[
\Weight{\greedy}{k} - \Weight{\antigreedy}{k} = \Root{\greedy}{k}.
\]
\end{lemm}

Observe that, as we have seen in Equation~\eqref{eq:rootbijection} on page~\pageref{eq:rootbijection}, this is also closely related to the bijection relating cluster algebras and subword complexes.

\begin{proof}[Proof of \Cref{lem:rootlemma}]
Starting with the greedy facet~$\greedy$, we flip the first position as long as we can without flipping into position~$k$ to obtain a facet~$I$. Observe that along these flips, the facet always consists of a consecutive sequence of the~$n$ simple reflections. We therefore obtain, up to commutations of consecutive commuting letters, that $I = \{k-n,\dots,k-1\}$ and
\[
\Weight{\greedy}{k} = \Weight{I}{k}.
\]
By the same argument, we flip the last position in $\antigreedy$ until we flip into position~$k$ to obtain a facet~$J$. Again up to commutations of consecutive commuting letters, we get that $J = \{k,\dots,k+n-1\}$ and
\[
\Weight{\antigreedy}{k} = \Weight{J}{k} = \Weight{J'}{k}.
\]
Here~$J'$ is the facet obtained from~$J$ by again flipping in~$J$ the last position~$n-1$ times so that, up to commutations, we have $J' = \{k-n+1,\dots,k\}$. The second equality is thus a direct consequence of \Cref{lem:weightequality} as~$k \in J \cap J'$. With these observations, we finally obtain for $\cwo{c} = \q_1\dots\q_N$ that
\begin{align*}
\Weight{\greedy}{k} - \Weight{\antigreedy}{k} &= \Weight{\{k-n,\dots,k-1\}}{k} - \Weight{\{k-n+1\dots,k\}}{k} \\
&= q_1\dots q_{k-1}(\omega_{q_k}) - q_1\dots q_{k}(\omega_{q_k}) \\
&= q_1\dots q_{k-1}(\omega_{q_k}) - q_1\dots q_{k-1}(\omega_{q_k} - \alpha_{q_k}) \\
&= q_1\dots q_{k-1}(\alpha_{q_k}) \\
&= \Root{\greedy}{k}. \qedhere
\end{align*}
\end{proof}
\setcounter{section}{3}
\section{Proof of Theorem~\ref{thm:main1}}\label{sec:proofmain1}

Before we recall the construction for the $F$-polynomials of type~$A_n$ to prove \Cref{thm:main1}, we set the needed notations. The Coxeter group~$W$ is the symmetric group $\Sym_{n+1}$ acting on $\RR^{n+1} \big/ \RR(1,\dots,1)$, whose simple system~$\sref$ is the set of simple transpositions $\sref=\{\tau_1,\dots,\tau_n\}$ for~$\tau_i = (i,i+1)$ interchanging $e_i$ and $e_{i+1}$. Thus, the Coxeter element~$c$ is given by the product of all simple transpositions in some order. The simple roots are moreover given by $\Delta = \set{e_i - e_{i+1}}{1 \leq i \leq n}$, the positive roots by $\Phiplus = \set{ e_i - e_{j+1}}{1 \leq i \leq j \leq n}$, and the fundamental weights by $\nabla = \set{e_1+\dots+e_i}{1 \leq i \leq n}$. We refer to \Cref{ex:subwordA2} on page~\pageref{ex:subwordA2} for these notations in type~$A_2$.

For consecutive simple transpositions, we write~$\tau_i < \tau_{i-1}$ if~$\tau_i$ appears to the left of~$\tau_{i-1}$ in~$c$, and~$\tau_i > \tau_{i-1}$ if~$\tau_i$ appears to the right of~$\tau_{i-1}$. We say that an element $\tau_{i_1} \dots \tau_{i_m}$ is a \Dfn{prefix of~$c$} if there is a reduced word~$\c$ for~$c$ beginning with~$\tau_{i_1},\dots,\tau_{i_m}$. If all $\tau_{i_1},\dots,\tau_{i_m}$ are inside the interval $\{\tau_i,\tau_{i+1},\dots,\tau_j\}$ for $i \leq j$, we moreover say that it is a \Dfn{prefix of~$c$ restricted to $\{\tau_i,\dots,\tau_j\}$} if the prefix property holds after removing all letters not in $\{\tau_i,\dots,\tau_j\}$ from~$c$. As an example, consider $c = \tau_1\tau_3\tau_2 = \tau_3\tau_1\tau_2 \in \Sym_{4}$. The prefixes of~$c$ are $-,\tau_1,\tau_3,\tau_1\tau_3 = \tau_3\tau_1,\tau_1\tau_3\tau_2$, and the prefixes of~$c$ restricted to $\{\tau_1,\tau_2\}$ are $-,\tau_1,\tau_1\tau_2$.

\subsection{\texorpdfstring{$F$}{F}-polynomials from \texorpdfstring{$T$}{T}-paths}\label{sec:fpolystpaths}

R.~Schiffler derived in~\cite{S2008} an explicit formula for the cluster variables of type~$A_n$ via \emph{T-paths}. These are certain paths on the diagonals of triangulations of a regular $(n+3)$-gon. G.~Musiker and R.~Schiffler then extended that description and obtained in~\cite{MS2010} an explicit formula for cluster variables in a similar fashion for cluster algebras with principal coefficients associated to unpunctured surfaces. Together with L.~Williams, they extended in~\cite{MSW2011} the results also to arbitrary surfaces, allowing punctures. In this section, we review that construction for type~$A_n$ to establish the needed notions to relate the description to the weight function in order to derive \Cref{thm:main1}. To present their results in a convenient way, we follow~\cite[Section~5]{MS2010} as they directly work with principal coefficients, except that we use slightly simplified notions of $T$-paths.

%\medskip

Let~$T$ be a triangulation of a regular $(n+3)$-gon, with boundary diagonals (or edges) labelled by $B_1,\dots,B_{n+3}$ and with proper diagonals labelled by $\tau_1,\dots,\tau_n$. An example can be found in \Cref{fig:A4trianexample}(a); we use $A,B,C,\dots$ instead of $B_1,B_2,B_3,\dots$ and $1,2,3,\dots$ instead of $\tau_1,\tau_2,\tau_3,\dots$ in examples for better readability.

Let $\gamma\notin T$ be another proper diagonal connecting non-adjacent vertices~$v_a$ and~$v_b$, oriented from~$v_a$ to~$v_b$. Denote the intersection points of~$\gamma$ with the diagonals in~$T$ along its orientation by $p_1,\dots,p_d$, and the corresponding diagonals in~$T$ by~$t_1,\dots,t_d$. Let $\gamma_{k}$ denote the segment of~$\gamma$ from point~$p_k$ to point~$p_{k+1}$, where we use $p_0=v_a$ and $p_{d+1}=v_b$. Each $\gamma_k$ lies in exactly one triangle~$\triangle_k$, and we orient the diagonal~$t_k$ in~$T$ by the orientation induced from the counterclockwise orientation of~$\triangle_k$. Note that if one considers the opposite path $\gamma^{-1}$ from~$v_b$ to~$v_a$, then the segment~$\gamma_i$ would become $\gamma_{d+1-i}$ and~$t_i$ would become $t_{d+1-i}$. Moreover, the induced orientations of all~$t_k$ would change.

A \Dfn{$T$-path}~$\zeta$ from~$v_a$ to~$v_b$ in~$T$ is a path~$\zeta=(\zeta_1,\dots,\zeta_{2d+1})$ in~$T$ where each $\zeta_i$ is either a diagonal in~$T$ or a boundary edge and which uses the oriented diagonals in the even positions, i.e.~$\zeta_{2k}=t_k$ for $1 \leq k \leq d$. Observe that such a $T$-path is uniquely determined by the directions in which the diagonals~$t_1,\dots,t_d$ in the even position are followed. If the direction of~$\zeta$ coincides along the diagonal~$t_k$ with the direction induced by the counterclockwise orientation of the triangle~$\triangle_k$, we write that~$\zeta$ travels~$t_k$ in \Dfn{positive direction}, and it travels~$t_k$ in \Dfn{negative direction} otherwise. It is not hard to see that there is always a unique $T$-path that travels all~$t_k$'s in positive direction. We call this path the \Dfn{greedy $T$-path}, and denote it by~$\zetag$. Similarly, we denote by~$\zetaag$ the \Dfn{antigreedy $T$-path} that travels all~$t_k$'s in negative direction. (Note that these appeared in~\cite{MS2010} as $\widetilde\alpha_{P_+}$ and $\widetilde\alpha_{P_-}$ in the paragraph before Theorem~5.1.) For instance, the greedy $T$-path in \Cref{fig:A4trianexample}(b) is $(F,2,3,3,3,4,B)$ and the antigreedy $T$-path is $(1,2,2,3,4,4,A)$.

We say that two $T$-paths~$\zeta$ and~$\zeta'$ are \Dfn{flipped} if~$\zeta$ and~$\zeta'$ only differ in two odd positions $2k-1$ and $2k+1$ for some~$k$. In other words,~$t_k$ is the unique diagonal which is traveled by~$\zeta$ and~$\zeta'$ in opposite directions, while all others are traveled in the same direction. We thus also say that~$t_k$ \Dfn{is flipped} between~$\zeta$ and~$\zeta'$. In \Cref{fig:A4trianexample}$\MK$(b), flipping~$\zeta_6 = t_3$ in the $T$-path $(F,2,3,3,3,4,B)$ yields the $T$-path $(F,2,3,3,C,4,A)$.

%\medskip

To a $T$-path~$\zeta=(\zeta_1,\dots,\zeta_{2d+1})$, one associates the monomial $m[\zeta]$ given by the product of variables $y_\ell$ such that~$\zeta_{2k} = t_k = \tau_\ell$ is traveled in positive direction. For instance, the greedy $T$-path~$\zetag = (F,2,3,3,3,4,B)$ yields the monomial $m[{\zetag}] = y_2y_3y_4$, while the antigreedy $T$-path~$\zetaag = (1,2,2,3,4,4,A)$ yields $m[{\zetaag}] = 1$. Moreover, all monomials obtained from $T$-paths for the diagonal~$\gamma$ in the example are given by
\[
\begin{array}{|C{4px}C{8px}C{4px}C{8px}C{4px}C{8px}C{8px}|c|}
\hline &&&\zeta&&& & m[\zeta] \\[2pt]
\hline &&&&&&&\\[-9pt] F&2^+&3&3^+&3&4^+&B & y_2y_3y_4 \\[2pt] F&2^+&3&3^+&C&4^-&A & y_2y_3 \\[2pt] 1&2^-&G&3^+&3&4^+&B & y_3y_4 \\[2pt] 1&2^-&G&3^+&C&4^-&A & y_3 \\[2pt] 1&2^-&2&3^-&4&4^-&A & 1 \\[2pt]
\hline
\end{array}
\]
where we labelled the even steps $\zeta_{2k} = t_k$ for $1 \leq k \leq d$ with a~``$+$'' if it is traveled in positive direction and with a~``$-$'' if it is traveled in negative direction. Also observe how $m[\zeta]$ changes under flips. If~$t_i$ is flipped between~$\zeta$ and~$\zeta'$ then $m[\zeta] = m[\zeta']\cdot y_\ell$ if~$t_i = \tau_\ell$ is traveled in~$\zeta$ in positive direction (and thus~$\zeta'$ in negative direction).

%\medskip

As we have noted above, the orientations of the~$t_k$'s depend on the orientation of~$\gamma$, while the monomial $m[\zeta]$ \emph{does not} depend on this orientation, but only on the two unordered endpoints $\{v_a,v_b\}$. This combinatorial model now provides a description of the $F$-polynomials for the cluster algebra where the initial datum is the fixed given triangulation~$T$ of a regular $(n+3)$-gon. It is well-known that $F$-polynomials for this cluster algebra are indexed by (unoriented) diagonals $\gamma \notin T$, see~\cite{MS2010}.

\begin{theo}[{\cite[Theorem~5.1]{MS2010}}]\label{thm:schiffler}
Let~$T$ be a triangulation of the regular $(n+3)$-gon, and let $\gamma \notin T$ with endpoints $\{v_a,v_b\}$. The $F$-polynomial $F_\gamma(\y)$ of~$\gamma$ is then given by
\[
F_\gamma(\y) = \sum m[\zeta],
\]
where the sum ranges over all $T$-paths~$\zeta$ from~$v_a$ to~$v_b$.
\end{theo}

%\medskip
Next, we recall how to associate a triangulation~$T_c$ of the regular $(n+3)$-gon to a Coxeter element~$c$ in type~$A_n$.
\begin{enumerate}
\item Pick a fixed vertex of the $(n+3)$-gon, labelled~$v_1$, and draw an edge connecting the two vertices adjacent to $v_1$. Label the new edge by $\tau_1$.
\item For each $i = 2,\dots,n$, label the vertex
\[
\begin{cases}
\text{clockwise} & \text{ if } \tau_i < \tau_{i-1} \\
\text{counterclockwise} & \text{ if } \tau_i > \tau_{i-1}
\end{cases}
\]
from~$v_{i-1}$ by~$v_i$, draw an edge connecting the two vertices adjacent to~$v_i$ and different from~$v_{i-1}$, and label the new edge~$\tau_i$.
\end{enumerate}

An example of type $A_3$ with $c = \tau_1\tau_3\tau_2$ is shown in \Cref{fig:132Tpaths}. Moreover, the triangulation in \Cref{fig:A4trianexample} corresponds to the Coxeter element $c = \tau_3\tau_2\tau_1\tau_4$ in type~$A_4$.

\begin{lemm}\label{lem:diagonalspositiveroots}
Let~$c$ be a Coxeter element in type~$A_n$, and let~$1 \leq i \leq j \leq n$. Then there is a unique diagonal~$\gamma \notin T_c$ that crosses exactly the diagonals labelled by $\tau_i,\tau_{i+1},\dots,\tau_j$, and every diagonal not in~$T_c$ can be obtained this way.
\end{lemm}

\begin{proof}
The triangulations that can be obtained (up to rotational symmetry) from a Coxeter element by the procedure are exactly the triangulations that do not have inner triangles, i.e. no triangles for which all three sides are proper diagonals. As the diagonals are labelled consecutively, the statement follows.
\end{proof}

We have the following corollary of the above \Cref{thm:schiffler}, which we will then use to deduce \Cref{thm:main1}.

\begin{coro}\label{cor:TpathAn}
Let~$c$ be a Coxeter element in type~$A_n$, let $\beta = e_i - e_{j+1}$ be a positive root, and let~$\gamma \notin T_c$ be the unique diagonal crossing exactly the diagonals labelled $\tau_i,\dots,\tau_j$ in~$T_c$. Let the endpoints of~$\gamma$ be $\{v_a,v_b\}$. The $F$-polynomials for the cluster algebra $\Alg(W,c)$ associated to~$\beta$ is then given by
\[
\Fpoly{\beta}{\y} = \sum m[\zeta],
\]
where the sum ranges over all $T$-paths~$\zeta$ from~$v_a$ to~$v_b$.
\end{coro}

\begin{proof}
This follows from the well known connection between $\Alg(W,c)$ and the triangulation~$T_c$ described above.
\end{proof}
We finally need the following proposition regarding possible flips in triangulations with respect to a Coxeter element~$c$.

\begin{prop}\label{prop:greedytpath}
Let~$T_c$ be the triangulation associated to a Coxeter element~$c$, and let~$\gamma\notin T_c$ be the unique diagonal oriented from $v_a$ to $v_b$ which crosses exactly the diagonals $\tau_i,\dots,\tau_j$ in this order. Then for any prefix $\tau_{i_1} \dots \tau_{i_m}$ of~$c$ restricted to $\{\tau_i,\dots,\tau_{j}\}$, one can flip the diagonals labelled~$\tau_{i_1},\dots,\tau_{i_m}$ in this order in the antigreedy $T_c$-path from~$v_a$ to~$v_b$. Moreover, every $T_c$-path from~$v_a$ to~$v_b$ is obtained this way for a unique prefix.
\end{prop}

\begin{proof}
We explicitly describe the four possible restrictions for directions in which $T_c$-paths can travel. To this end, consider the situation that~$t_{i-1}$ and~$t_{i}$ are oriented towards their shared vertex in~$T_c$, or, equivalently, that $\tau_{i} < \tau_{i-1}$ (see the path for $e_2-e_4$ in \Cref{fig:132Tpaths}). Then, any~$T_c$-path $\zeta$ from~$v_a$ to~$v_b$
\begin{itemize}
\item that travels the diagonal~$\zeta_{2i-2} = t_{i-1}$ in \emph{positive} direction must also travel $\zeta_{2i} = t_{i}$ in \emph{positive} direction, and
\item that travels the diagonal~$\zeta_{2i} = t_i$ in \emph{negative} direction must also travel $\zeta_{2i-2} = t_{i-1}$ in \emph{negative} direction.
\end{itemize} The situation where~$t_{i-1}$ and~$t_{i}$ are oriented away from their shared vertex in~$T_c$, or, equivalently, that $\tau_i > \tau_{i-1}$ is the same with the roles of positive and negative direction interchanged (see the path for $e_1-e_3$ in \Cref{fig:132Tpaths}).

Clearly, these are the only restrictions on $T_c$-paths. This means that a given sequence of orientations of the $t_i$'s corresponds to a $T_c$-path if and only if these restrictions are satisfied. It then directly follows that every $T_c$-path is uniquely obtained from the antigreedy $T_c$-path~$\zetaag$ (which travels all the diagonals in negative direction) by flipping diagonals labelled $\tau_{i_1},\dots,\tau_{i_m}$ in this order for a prefix $\tau_{i_1} \dots \tau_{i_m}$ of~$c$ restricted to $\{\tau_i,\dots,\tau_{j}\}$, as desired. Observe here that if one considers two different words for the same prefix, then both sequences of flips yield the same $T_c$-path, as expected.
\end{proof}

We refer to \Cref{fig:132Tpaths} for several examples, and also to \Cref{ex:A3two} below for two concrete computations.

\begin{coro}\label{cor:TpathAn2}
In the situation of \Cref{cor:TpathAn}, we have that
\[
\Fpoly{\beta}{\y} = \sum y_{i_1}\dots y_{i_m}
\]
where the sum ranges over all prefixes $\tau_{i_1} \dots \tau_{i_m}$ of~$c$ restricted to $\{\tau_i,\dots,\tau_{j}\}$.
\end{coro}

\begin{proof}
\Cref{prop:greedytpath} implies that $\Fpoly{\beta}{\y} = \sum m[\zeta_{i_1,\dots,i_m}]$ where the sum ranges over all prefixes of~$c$ restricted to $\{\tau_i,\dots,\tau_j\}$. Moreover, the resulting $T$-path $\zeta_{i_1,\dots,i_m}$ is obtained from the antigreedy $T$-path $\zetaag$ by flipping all the diagonals $\tau_{i_1},\dots,\tau_{i_m}$ in this order. The definition of $m[\zeta]$ thus implies that $m[\zeta_{i_1,\dots,i_m}] = y_{i_1}\dots y_{i_m}$, as desired.
\end{proof}
\subsection{\texorpdfstring{$F$}{F}-polynomials from subword complexes}

We now use \Cref{cor:TpathAn2} to obtain the $F$-polynomials from the weight vectors of $\clusterComplex$ by providing the analogous property of the weight vectors in \Cref{cor:intervalequality}. We start with the following explicit description of all different weights $\Weight{I}{k}$ that occur for the various facets for a fixed position~$k$. For $1 \leq i \leq j \leq n$, consider the positive root $e_i - e_{j+1}$. We have seen in \Cref{lem:lemma33} that there is a unique~$k \in \{n+1,\dots,n+N\}$ such that $\Root{\greedy}{k} = e_i - e_{j+1}$, and we have then seen in \Cref{lem:rootlemma} that
\begin{align}
\Weight{\greedy}{k} - \Weight{\antigreedy}{k} = \Root{\greedy}{k} = e_i - e_{j+1}. \label{eq:weighttypeA}
\end{align}

This yields the following lemma.

\begin{lemm}\label{lem:greedyantigreedyweight}
We have
\begin{align}
\Weight{\greedy}{k} &= (\epsilon_1,\dots,\epsilon_{i-1},\ 1,\ \epsilon_{i+1},\dots,\epsilon_{j},\ 0,\ \epsilon_{j+2},\dots,\epsilon_{n+1}) \label{eq:greedyAn}
\\
\Weight{\antigreedy}{k} &= (\epsilon_1,\dots,\epsilon_{i-1},\ 0,\ \epsilon_{i+1},\dots,\epsilon_{j},\ 1,\ \epsilon_{j+2},\dots,\epsilon_{n+1}) \label{eq:antigreedyAn}
\end{align}
for fixed $\epsilon_i \in \{0,1\}$ with $i \in \{1,\dots,n+1\}\setminus\{i,j+1\}$.
\end{lemm}

\begin{proof}
This follows from~\eqref{eq:weighttypeA} as all $W$-orbits of fundamental weights in type~$A_n$ consist of $(0,1)$-vectors.
\end{proof}

Using this observation, one explicitly obtains these weights for a given Coxeter element~$c$ as described next.

\begin{prop}\label{prop:weightsAn}
In the situation of \Cref{lem:greedyantigreedyweight}, we have the following properties of the $\epsilon_i$'s (where the cases in~\eqref{eq:innerepsilon} are only considered either if $i>1$, or if $j<n$, respectively):
\begin{enumerate}%[(i)]
\item $\epsilon_1 = \dots = \epsilon_{i-1}$ and $\epsilon_{j+2} = \dots = \epsilon_{n+1}$;
\item\label{eq:innerepsilon}
$\epsilon_{i-1} =
\begin{cases}
0 &\text{ if } \tau_{i-1} < \tau_{i } \\
1 &\text{ if } \tau_{i-1} > \tau_{i }
\end{cases}$\quad and\quad $\epsilon_{j+2} =
\begin{cases}
0 &\text{ if }\tau_{j+1} < \tau_{j+2} \\
1 &\text{ if }\tau_{j+1} > \tau_{j+2};
\end{cases}$
\item\label{it:innerepsilon}
for $i<\ell\leq j$, $\epsilon_\ell =
\begin{cases}
1 \text{ if }\tau_{\ell-1} < \tau_\ell \\
0 \text{ if }\tau_{\ell-1} > \tau_\ell.
\end{cases}$
\end{enumerate}
\end{prop}

\begin{proof}
One can flip the letters in the initial copy of $\c$ inside $\cw{c}$ from right to left to obtain a sequence of increasing flips from~$\greedy$ to~$\antigreedy$. (We indeed obtain exactly the shortest sequences of increasing flips from the greedy to the antigreedy facet.) As the root configuration of~$\greedy$ is given by all simple roots $\{e_1-e_2,\dots,e_n-e_{n+1}\}$, \Cref{lem:lemma33}$\MK$\eqref{it:rootupdate} implies that along the above sequence of flips from~$\greedy$ to~$\antigreedy$, every pattern~$(\epsilon_\ell,\epsilon_{\ell+1})$ of consecutive indices of $\Weight{\greedy}{k}$ is updated exactly once (where we set $\epsilon_i = 1$ and $\epsilon_{j+1} = 0$ as in~\eqref{eq:greedyAn}. Moreover, \Cref{lem:positiveweightshift} implies that along this procedure, either $\epsilon_\ell = \epsilon_{\ell+1}$ and the application of~$\tau_\ell$ does not change the weight, or $(\epsilon_\ell,\epsilon_{\ell+1}) = (1,0)$, and this application moves the $1$ in position~$\ell$ into position~$\ell+1$. As such a move is not reversible again by \Cref{lem:positiveweightshift}, we directly obtain the second property. The first property is obtained with the additional observation that those entries coincide in~$\greedy$ and in~$\antigreedy$. The last property finally follows with the observation that every $\tau_\ell$ for $i \leq \ell \leq j$ must indeed move a~$1$ one position to the right to obtain~$\antigreedy$ from~$\greedy$ this way.
\end{proof}
To state the main observation towards the proof of \Cref{thm:main1}, define the following set of weights as an ``interval'' in the weights,
\[
\IntervalWeight{k} = \set{\omega \in W(\nabla)}{\Weight{\greedy}{k} - \omega, \omega - \Weight{\antigreedy}{k} \in L^+}.
\]

Using this notion, we deduce the following theorem from \Cref{prop:weightsAn}, for which we need one additional observation. We may flip from the greedy to the antigreedy facet in $\clusterComplex$ in~$n$ flips. For a fixed expression $c = s_1\dots s_n$, these flips are exactly given by sequences of flips in positions $i_n,\dots,i_1$ such that $c = s_{i_1}\dots s_{i_n}$. We denote by $\mathcal{I}$ the set of facets of $\clusterComplex$ that lie on such a shortest sequence of increasing flips and remark that $\mathcal{I}$ is given by construction by all facets~$I$ that are obtained from $\greedy$ by flipping \emph{suffixes} of the Coxeter element~$c$, compare \Cref{ex:shortestpaths}.

\begin{theo}\label{cor:intervalequality}
We then have for $k \in \{1,\dots,n+N\}$ that
\begin{equation*}
\set{\Weight{I}{k}}{I \text{ facet of }\clusterComplex} = \set{\Weight{I}{k}}{I \in \mathcal{I}}.
\end{equation*}
\end{theo}

\begin{proof}
The first inclusion in
\begin{align*}
\set{\Weight{I}{k}}{I \in \mathcal{I}} &\subseteq \set{\Weight{I}{k}}{I \text{ facet of }\clusterComplex} \\
&\subseteq \IntervalWeight{k}.
\end{align*}
is trivial, while the second is a direct consequence of \Cref{lem:positiveweightshift} as every facet lies on a (not necessarily shortest) sequence of increasing flips from the greedy to the antigreedy facet. On the other hand, \Cref{prop:weightsAn} implies that $\set{\Weight{I}{k}}{I \in \mathcal{I}} = \IntervalWeight{k}$ as follows. For $\Root{\greedy}{k} = e_i - e_{j+1}$, we have that the interval $\IntervalWeight{k}$ is given by starting with $\Weight{\greedy}{k}$ and replacing consecutive pattern $(\epsilon_\ell,\epsilon_{\ell+1}) = (1,0)$ by $(0,1)$ in all possible ways so that every consecutive pattern $(\epsilon_\ell,\epsilon_{\ell+1})$ with $i \leq \ell \leq j$ is modified exactly once. Now observe that~\Cref{prop:weightsAn}$\MK$\eqref{it:innerepsilon} gives that $(\epsilon_\ell,\epsilon_{\ell+1}) = (1,0)$ in $\Weight{\greedy}{k}$ if and only if $\tau_{\ell-1} < \tau_\ell > \tau_{\ell+1}$ for $i < \ell < j$ and $\tau_\ell > \tau_{\ell+1}$ for $\ell = i$ and $\tau_{\ell-1} < \tau_\ell$ for $\ell = j$. This means that the possible ways to replace consecutive patterns $(\epsilon_\ell,\epsilon_{\ell+1}) = (1,0)$ by $(0,1)$ in $\Weight{\greedy}{k}$ is given by doing this replacements at all indices corresponding to \emph{suffixes} of the Coxeter element~$c$ when restricted to $\tau_i,\dots,\tau_j$.
\end{proof}
\begin{proof}[Proof of \Cref{thm:main1}]
Let $\Alg(W,c)$ be the cluster algebra of type~$A_n$ for a given Coxeter element~$c$, and let $\beta = e_i - e_{j+1}$ be a positive root. By \Cref{cor:TpathAn2}, we have that
\begin{equation*}\label{eq:Fpolysum}
\Fpoly{\beta}{\y} = \sum y_{i_1}\dots y_{i_m}
\end{equation*}
where the sum ranges over all prefixes $\tau_{i_1} \dots \tau_{i_m}$ of~$c$ restricted to $\{\tau_i,\dots,\tau_{j}\}$. Let~$k$ be the unique index such that $\Root{\greedy}{k} = \beta$. We aim to show that this sum is also given by $\sum \y^\gamma$ where the sum ranges over all $\gamma$ in
\[
\bigset{ \Weight{I}{k} - \Weight{\antigreedy}{k}}{ I \in \mathcal{I} }
\]
expressed in the root basis~$\Delta$. Together with \Cref{cor:intervalequality}, this implies \Cref{conj:newtonpolytope} for $\Alg(W,c)$. \Cref{conj:latticepoints} is then trivial as all exponent vectors of monomials in $\Fpoly{\beta}{\y}$ are $(0,1)$-vectors in type~$A$.

%\medskip

By construction, the set~$\mathcal{I}$ in \Cref{cor:intervalequality} is obtained from the greedy facet $\greedy$ by flipping all possible sequences of \emph{suffixes} of the Coxeter element~$c$. Let $i_1,\dots,i_n$ be such that $c = \tau_{i_1} \dots \tau_{i_n}$, and let~$I_m$ be the facet obtained from~$\greedy$ by flipping the letters $i_n,i_{n-1},\dots,i_m$ in this order. Then flipping the letter $i_n$ in $\greedy$ yields $I_n$, flipping the letter $i_m$ in $I_{m+1}$ yields $I_m$ and $I_1 = \antigreedy$. Moreover,
\begin{equation}\label{eq:rootupdateshortestpath}
\Root{\greedy}{i_m} = \Root{I_{m+1}}{i_m} = \alpha_{i_m}
\end{equation}
because~$\greedy$ and~$I_{m+1}$ coincide until position~$i_m$. This means that $\Weight{I_{m+1}}{k} - \Weight{I_m}{k}$ is either~$0$ or $\alpha_{i_m}$. Since $\Weight{\greedy}{k} - \Weight{\antigreedy}{k} = \Root{\greedy}{k}$, this implies that $\Weight{I_{m}}{k} - \Weight{\antigreedy}{k}$ is obtained from $\Root{\greedy}{k}$ by replacing all $1$'s by $0$'s in positions $\{i_m,\dots,i_n\}$. This is because every position that is $1$ in $\Root{\greedy}{k}$ must become $0$ within the sequence $\greedy,I_n,\dots,I_1=\antigreedy$ of increasing flips, and~\eqref{eq:rootupdateshortestpath} means that the only flip that may modify a position~$1 \leq a \leq n$ is the flip of the index~$i_m$ with $i_m = a$.

In total, this gives that $\sum \y^\gamma$ ranging over all $\gamma \in \set{ \Weight{I}{k} - \Weight{\antigreedy}{k}}{ I \in \mathcal{I} }$ equals $\sum \y^\beta / (y_{i_1}\dots y_{i_m})$ ranging over all suffixes of~$c$ restricted to $\{\tau_i,\dots,\tau_j\}$. The later is clearly equal to the desired sum expression for $\Fpoly{\beta}{\y}$.
\end{proof}
\begin{thebibliography}{28}
\bibitem[10]{Hum1990} James E. Humphreys, Reflection groups and Coxeter groups, vol. 29, Cambridge University Press, 1990.
\bibitem[6]{FZ20022} Sergey Fomin and Andrei Zelevinsky, Cluster algebras I: foundations, J. Am. Math. Soc. 15 (2002), no. 2, 497–529.
\bibitem[7]{FZ2003} Sergey Fomin and Andrei Zelevinsky, Cluster algebras II: finite type classification, Invent. Math. 154 (2003), no. 1, 63–121.
\bibitem[8]{FZ2007} Sergey Fomin and Andrei Zelevinsky, Cluster algebras IV: coefficients, Compos. Math. 143 (2007), no. 1, 112–164.
\bibitem[28]{YZ2008} Shih-Wei Yang and Andrei Zelevinsky, Cluster algebras of finite type via Coxeter elements and principal minors, Transform. Groups 13 (2008), no. 3-4, 855–895.
\bibitem[11]{KM2004} Allen Knutson and Ezra Miller, Subword complexes in Coxeter groups, Adv. Math. 184 (2004), no. 1, 161–176.
\bibitem[3]{CLS2014} Cesar Ceballos, Jean-Philippe Labbé, and Christian Stump, Subword complexes, cluster complexes, and generalized multi-associahedra, J. Algebr. Comb. 39 (2014), no. 1, 17–51.
\bibitem[22]{R2007} Nathan Reading, Sortable elements and Cambrian lattices, Algebra Univers. 56 (2007), no. 3-4, 411–437.
\bibitem[19]{PS2015} Vincent Pilaud and Christian Stump, Brick polytopes of spherical subword complexes and generalized associahedra, Adv. Math. 276 (2015), 1–61.
\bibitem[25]{S2008} Ralf Schiffler, A cluster expansion formula (An case), Electron. J. Comb. 15 (2008), R64 (9 pages).
\bibitem[16]{MS2010} Gregg Musiker and Ralf Schiffler, Cluster expansion formulas and perfect matchings, J. Algebr. Comb. 32 (2010), no. 2, 187–209.
\bibitem[17]{MSW2011} Gregg Musiker, Ralf Schiffler, and Lauren Williams, Positivity for cluster algebras from surfaces, Adv. Math. 227 (2011), no. 6, 2241–2308.
\end{thebibliography}
\end{document}
